\subsection*{Performance by data type}

Data type separated model performance sharply (Figure~\ref{fig:data-types}). Mean pass rates were lowest on the less common omics readouts, such as TCR-seq, metabolomics and proteomics (14.1\%, 10 evaluations), and on clinical and biofluid data (17.6\%, 7). The best performance was observed for EEG, ECoG and sleep scoring (25.9\%, 10), histology and imaging (25.6\%, 12), behavior and pose tracking (24.8\%, 11) and bulk RNA-seq (24.6\%, 18). The two most common data types, single-cell RNA-seq (16.3\%, 29) and extracellular electrophysiology or BCI decoding (19.6\%, 29), sat below the mean over models of 20.1\%. Opus 5.5 passed 50.0\% of bulk RNA-seq evaluations but none of those using patch clamp, multi-electrode arrays or calcium imaging, whereas GPT-6 Astra passed 35.7\% of the latter and only 11.9\% of clinical and biofluid evaluations. Across the 21 models, the gap between a model's best and worst pass rate averaged 28 points, indicating that the biological context carried by the data type strongly shapes agent performance.
